\clearpage
\VerdierPhysicalPageAnchor{0177}%
\begingroup\small\setlength{\parskip}{0pt}\setlength{\emergencystretch}{1.5em}
\setlength{\abovedisplayskip}{4pt plus 2pt minus 2pt}%
\setlength{\belowdisplayskip}{4pt plus 2pt minus 2pt}%

\VerdierSemanticAnchor{ch3d-1-2-1-truncations-cont}
\noindent(respectively,
\[
\begin{aligned}
i^q&=\operatorname{id},&&q<n-1,\\
i^{n-1}&\colon\operatorname{Ker}(d_X^{n-1})\longrightarrow X^{n-1}
  &&\text{(canonical morphism)},\\
i^q&=0,&&q\geq n
\end{aligned}
\]
).

Notice that the morphism $p$ (respectively, $i$) induces an isomorphism on
the cohomology objects in degree $j\geq n$ (respectively, $j<n$), and that
$H^j(X(n,+\infty))=0$ for $j<n$ (respectively,
$H^j(X(-\infty,n))=0$ for $j\geq n$). All the assertions of Proposition
(1.2.1) now follow immediately from this construction. We leave the
detailed proof to the reader.

\VerdierSemanticAnchor{ch3d-1-2-2}
\medskip\noindent\textbf{Definition 1.2.2.}\quad
The triangulated category
\[
D(\mathcal A)=K(\mathcal A)/\operatorname{Ac}(\mathcal A)
=K(\mathcal A)[\operatorname{Qis}(\mathcal A)^{-1}]
\]
is called the \emph{derived category of $\mathcal A$} and is denoted by
$D(\mathcal A)$ (Chap. II, 2.2.6 and 2.2.10).

It follows from Theorem (2.2.6) that the canonical cohomological functor
$H\colon K(\mathcal A)\longrightarrow\mathcal A$ factors uniquely through
a cohomological functor, still denoted by

\VerdierSemanticAnchor{ch3d-1-2-2-cohomology}
\begin{equation}
H\colon D(\mathcal A)\longrightarrow\mathcal A,
\tag{1.2.2.1}
\end{equation}
and likewise called the \emph{canonical cohomological functor}. The
canonical quotient functor is denoted by
$Q\colon K(\mathcal A)\longrightarrow D(\mathcal A)$.

\VerdierSemanticAnchor{ch3d-1-2-3}
\medskip\noindent\textbf{Theorem 1.2.3.}\quad
\textit{a) The natural functors appearing in the diagram:}

\VerdierSemanticAnchor{ch3d-1-2-3-a-diagram}
\begin{center}
  \resizebox{0.80\textwidth}{!}{%
    $\begin{tikzcd}[ampersand replacement=\&,column sep=4.5em,row sep=3.0em]
      \& K^{,+}(\mathcal A)/\operatorname{Ac}(\mathcal A)
        \arrow[dr] \& \\
      K^{,b}(\mathcal A)/\operatorname{Ac}(\mathcal A)
        \arrow[ur] \arrow[dr]
        \& \& D(\mathcal A) \\
      \& K^{,-}(\mathcal A)/\operatorname{Ac}(\mathcal A)
        \arrow[ur] \&
    \end{tikzcd}$%
  }
\end{center}
\endgroup
